\usepackage{geometry}
\geometry{top=2.5cm, bottom=2.5cm, left=2.2cm, right=2.2cm}

% --- การตั้งค่าภาษาไทยและฟอนต์ ---
\usepackage{fontspec}
\setmainfont{TH Sarabun New} % สำหรับโปรแกรมในเครื่องคอมพิวเตอร์ที่มีฟอนต์นี้
% หากใช้ Overleaf แนะนำให้ดาวน์โหลดฟอนต์ TH Sarabun New มาใส่ในโปรเจกต์ หรือเปลี่ยนเป็น \setmainfont{Sarabun}

\usepackage{amsmath,amssymb}
\usepackage{graphicx}
\usepackage{fancyhdr}
\usepackage{lastpage}
\usepackage{tabularx}

% เพิ่มระยะห่างระหว่างบรรทัดเล็กน้อย เพื่อไม่ให้สระและวรรณยุกต์ภาษาไทยซ้อนทับกัน
\linespread{1.25}

% --- การตั้งค่าหัวกระดาษ (Header) แสดงเลขหน้า "หน้า X/Y" ---
\pagestyle{fancy}
\fancyhf{} % ล้างค่าเริ่มต้นทั้งหมด
\fancyhead[R]{\fontsize{12pt}{12pt}\selectfont หน้า \thepage} % จัดให้ชิดขวาตามข้อสอบจริง
\renewcommand{\headrulewidth}{0pt} % เอาเส้นคั่นแนวนอนออก
\renewcommand{\footrulewidth}{0pt}

% --- ตัวนับเลขข้อแบบรันต่อเนื่อง (Continuous Question Counter) ---
\newcounter{qnum}
\newcommand{\question}{\refstepcounter{qnum}\par\medskip\noindent\textbf{\theqnum. }}

% --- คำสั่งสำหรับสร้างหัวข้อแต่ละตอน ---
\newcommand{\examsection}[1]{%
  \vspace{0.6cm}
  \noindent\textbf{\large #1}
  \vspace{0.3cm}\par
}

% --- ระบบจัดรูปแบบช้อยส์ (ก. ข. ค. ง.) ---
% 1. ช้อยส์แบบแนวนอนเรียงแถวเดียว (เหมาะสำหรับตัวเลขสั้นๆ)
\newcommand{\choiceshorizontal}[4]{%
  \nopagebreak
  \noindent\makebox[0.22\textwidth][l]{ก. #1}%
  \makebox[0.22\textwidth][l]{ข. #2}%
  \makebox[0.22\textwidth][l]{ค. #3}%
  \makebox[0.22\textwidth][l]{ง. #4}\par\medskip
}

% 2. ช้อยส์แบบตาราง 2x2 (เรียง ก. ข. อยู่บน ค. ง. อยู่ล่าง เหมาะสำหรับข้อความสั้น-ปานกลาง)
\newcommand{\choicesgrid}[4]{%
  \nopagebreak
  \noindent\begin{tabular}{@{}p{0.45\textwidth}p{0.45\textwidth}}
    ก. #1 & ข. #2 \\
    ค. #3 & ง. #4
  \end{tabular}\par\medskip
}

% 3. ช้อยส์แบบเรียงดิ่งลงมาแถวละข้อ (เหมาะสำหรับช้อยส์ที่เป็นประโยคยาวๆ)
\newcommand{\choicesvertical}[4]{%
  \nopagebreak
  \noindent\hspace*{1em} ก. #1 \\
  \noindent\hspace*{1em} ข. #2 \\
  \noindent\hspace*{1em} ค. #3 \\
  \noindent\hspace*{1em} ง. #4 \par\medskip
}
